\section{Estructura y parser del \gls{BT}}
\label{sec:traduccionBT}

HABLAR PRIMERO DEL MODELO

Para poder generar y traducir después al \gls{motor} los \glspl{BT} es necesario especificar qué estructura deben de seguir y qué son capaces de hacer.
Este trabajo se ha realizado en el motor Unreal Engine 5(METER CITA O ENTRADA DE GLOSARIO SUPONGO), por lo que las capacidades de los \glspl{BT} que se buscan crear deben de ajustarse con las posibilidades que ofrece este \gls{motor}.

\subsection{\glspl{BT} de Unreal Engine.}
\label{sec:UEBT}
NO SÉ SI PONERLO AQUÍ O EN EL ESTADO DEL ARTE O DONDE PERO POR AHORA LO DEJO AQUÍ.
TAMPOCO SÉ SI DEBERÍA EXPLICAR CONCEPTOS COMO ACTORES, COMPONENTES O COSAS ASÍ DE UNREAL, Y SI ES QUE SÍ DÓNDE

Los \glspl{BT} se evalúan en cada tick por el personaje(GLOSARIO) que lo posee.
Las evaluaciones se realizan en los nodos, que pueden ser de diferentes familias.

Los pimeros nodos que nos encontramos son los intermedios, llamados ``Composite''.
Estos nodos no pueden ser terminales y sirven como control de flujo para decidir cómo se van a ejecutar sus nodos hijos.
Dependiendo de cómo se comportan en cuanto a la elección de ejecución de sus hijos se clasifican en tres:
\begin{enumerate}
	\renewcommand{\theenumi}{\alph{enumi}}
	\item Selector: Ejecuta solo el primer nodo hijo que pueda.
	\item Sequence: Ejecuta todos sus hijos en órden hasta terminar o hasta que un nodo devuelva fallo.
	\item Simple Parallel: Tiene dos ramas. La primera ejecuta una tarea (rama de tarea principal), mientras que la segunda ejecuta a la vez una actividad. 
\end{enumerate}

Los siguientes nodos que nos podemos encontrar son los decoradores.
Actúan como nodos condicionales, acoplándose a un nodo y condicionando si ese nodo puede ejecutarse o no.

Finalmente nos encontramos con los nodos terminales (``Tasks''), donde se encuentran las acciones que deben de ser ejecutadas.
Estos nodos pueden necesitar variables de entrada/salida, guardadas en una estructura llamada ``Blackboard''(GLOSARIO).

Esta Blackboard es una estructura de datos que almacena información en forma de clave-valor, comportándose como memoria compartida entre el \gls{BT} y el componente de IA.
Las claves contenidas pueden usarse para comprobar condiciones de los decoradores o para almacenar datos necesarios para las tareas.

(NO SÉ SI PONER CAPTURA DE LA PÁGINA DE UNREAL DE EJEMPLO DE BT)

BT Y EL AICONTROLLER



\subsection{Estructura del \gls{BT}}
JSON SCHEMA
Teniendo en cuenta la estructura que tiene Unreal Engine, se decidió diseñar un esquema de cómo deberían generarse los \glspl{BT} por el \gls{LLM} para poder ser validados.
El formato elegido para la creación de \glspl{BT} ha sido el formato JSON debido a su fácil legibilidad, por lo que el esquema de validación se ha creado siguiendo la especificación de JSON Schema 2020-12


\subsection{Parser de \glspl{BT}}
